% Biên dịch bằng XeLaTeX (cần cho tiếng Việt có dấu).
%   xelatex nxcred-overview.tex
% Hoặc tải cả thư mục docs/ lên Overleaf rồi chọn compiler XeLaTeX.
\documentclass[a4paper,11pt]{article}

\usepackage{fontspec}
\setmainfont{TeX Gyre Termes}
\setsansfont{TeX Gyre Heros}
\setmonofont{TeX Gyre Cursor}

\usepackage[margin=2.6cm]{geometry}
\usepackage{tikz}
\usetikzlibrary{arrows.meta,positioning,calc,fit,backgrounds}
\usepackage{enumitem}
\usepackage{booktabs}
\usepackage{longtable}
\usepackage{xcolor}
\usepackage[hidelinks]{hyperref}
\usepackage{titlesec}

\definecolor{ink}{HTML}{16171A}
\definecolor{muted}{HTML}{6E7079}
\definecolor{accent}{HTML}{2F5FE0}
\definecolor{good}{HTML}{17795E}
\definecolor{line}{HTML}{D8D8D2}
\definecolor{soft}{HTML}{F5F5F1}

\hypersetup{colorlinks=true,linkcolor=accent,urlcolor=accent}

\titleformat{\section}{\sffamily\LARGE\bfseries\color{ink}}{\thesection.}{0.5em}{}
\titleformat{\subsection}{\sffamily\large\bfseries\color{ink}}{\thesubsection}{0.5em}{}

\newcommand{\lead}[1]{{\color{muted}\large #1}\par\medskip}

% Một dòng mục lục kiểu ethereum.org: "Tên mục -- mô tả ngắn"
\newcommand{\topicline}[2]{\item \textbf{#1} \, --- \, {\color{muted}#2}}

\tikzset{
  blk/.style={draw=line,fill=white,rounded corners=4pt,inner sep=7pt,
              font=\sffamily\small,align=center,minimum height=10mm},
  actor/.style={blk,fill=soft,font=\sffamily\small\bfseries,minimum width=26mm},
  chain/.style={blk,fill=soft,draw=accent,font=\sffamily\small\bfseries},
  step/.style={-{Stealth[length=2.4mm]},draw=ink,thick,
               font=\sffamily\scriptsize,align=center},
  note/.style={font=\sffamily\scriptsize\itshape,color=muted,align=center},
}

\begin{document}

\begin{center}
{\sffamily\Huge\bfseries NX Cred}\\[4pt]
{\sffamily\Large\color{muted} Tài liệu giới thiệu hệ thống định danh số}\\[10pt]
{\color{muted}\small Bản demo --- cập nhật \today}
\end{center}

\vspace{6mm}

\lead{Tài liệu này giới thiệu NX Cred: một hệ thống cấp và xác minh giấy tờ định danh mà
bên xác minh có thể tin chắc bạn là bạn, nhưng không nhìn thấy dữ liệu gốc của bạn.
Tài liệu đi từ khái niệm tới kiến trúc rồi tới các luồng hoạt động, không đòi hỏi người
đọc biết trước về mật mã học.}

\subsection*{Nên bắt đầu từ đâu}

Mỗi người đến với tài liệu này vì một mục đích khác nhau. Vài lối vào thường gặp:

\begin{itemize}[leftmargin=1.2em,itemsep=4pt]
  \item \textbf{Muốn hiểu ý tưởng trong 5 phút.} Đọc mục \ref{sec:zkp} và \ref{sec:nxcred},
        rồi xem sơ đồ ở mục \ref{sec:flow-hv}. Bỏ qua phần còn lại.
  \item \textbf{Muốn biết hệ thống gồm những gì.} Đọc mục \ref{sec:topology} rồi tới
        phần \textit{NX stack}.
  \item \textbf{Muốn thẩm định về mặt kỹ thuật.} Đọc phần \textit{Nền tảng thiết kế},
        sau đó tải tài liệu lõi mật mã ở mục \ref{sec:crypto}.
  \item \textbf{Muốn đọc tuần tự từ đầu tới cuối.} Các mục dưới đây đã xếp theo đúng
        thứ tự đó.
\end{itemize}

\vspace{2mm}
\hrule
\vspace{6mm}

\part*{\sffamily Overview}
\addcontentsline{toc}{part}{Overview}

\section{Topology}
\label{sec:topology}

NX Cred có bốn thành phần. Ba bên tham gia và một sổ cái công khai.

\begin{center}
\begin{tikzpicture}[node distance=16mm]
  \node[chain] (chain) {Chain\\[2pt]\footnotesize\mdseries sổ cái công khai};
  \node[actor,below left=20mm and 26mm of chain] (issuer) {Issuer\\[2pt]\footnotesize\mdseries cơ quan cấp};
  \node[actor,below=34mm of chain] (holder) {Holder\\[2pt]\footnotesize\mdseries ví của bạn};
  \node[actor,below right=20mm and 26mm of chain] (verifier) {Verifier\\[2pt]\footnotesize\mdseries bên xác minh};

  \draw[step] (issuer) -- node[left,xshift=-1mm]{đăng khoá\\công khai} (chain);
  \draw[step] (chain) -- node[right,xshift=1mm]{đọc khoá\\công khai} (verifier);
  \draw[step] (issuer) -- node[below,sloped]{cấp thẻ} (holder);
  \draw[step] (holder) -- node[below,sloped]{trình bằng chứng} (verifier);

  \node[note,below=3mm of holder] {Dữ liệu gốc không bao giờ rời khỏi ví};
\end{tikzpicture}
\end{center}

Hai điều đáng chú ý trong sơ đồ này, vì chúng quyết định mọi thứ phía sau:

\begin{itemize}[leftmargin=1.2em,itemsep=4pt]
  \item \textbf{Issuer và Verifier không bao giờ nói chuyện trực tiếp.} Verifier không
        cần gọi điện hỏi cơ quan cấp xem thẻ này có thật không. Nó chỉ cần đọc khoá công
        khai đã đăng trên chain. Cơ quan cấp do đó \emph{không biết} bạn đang dùng thẻ ở đâu.
  \item \textbf{Giao dịch xác minh không đi qua chain.} Việc trình bằng chứng diễn ra
        trực tiếp giữa ví và bên xác minh. Chain chỉ giữ khoá công khai và khuôn mẫu
        giấy tờ, không giữ lịch sử ai xác minh cái gì.
\end{itemize}

\section{Zero-knowledge proof là gì}
\label{sec:zkp}

Zero-knowledge proof --- bằng chứng không tiết lộ --- là cách chứng minh một điều là đúng
mà không nói ra vì sao nó đúng.

Một ví dụ quen thuộc: bạn muốn chứng minh mình đủ 18 tuổi để vào một quán bar. Cách thông
thường là đưa căn cước ra. Nhưng tấm căn cước đó nói cho nhân viên biết luôn cả tên bạn,
ngày sinh chính xác, quê quán và địa chỉ nhà --- toàn bộ những thứ họ không cần và không
nên biết. Bạn trả lời một câu hỏi bằng cách tiết lộ mười câu trả lời.

Bằng chứng không tiết lộ cho phép trả lời đúng một câu hỏi được hỏi. Ba tính chất của nó:

\begin{itemize}[leftmargin=1.2em,itemsep=4pt]
  \item \textbf{Đầy đủ} --- điều đúng thì luôn chứng minh được.
  \item \textbf{Chắc chắn} --- điều sai thì không cách nào chứng minh được, kể cả khi
        người chứng minh có máy tính mạnh và cố tình gian lận.
  \item \textbf{Không tiết lộ} --- bên kiểm tra kết thúc phiên mà không biết thêm bất cứ
        điều gì ngoài đúng một chữ ``đúng''.
\end{itemize}

Điều khiến nó không phải là ma thuật: bên chứng minh phải trả lời một câu hỏi ngẫu nhiên
mà họ \emph{không thể đoán trước}. Đoán mò thì xác suất qua được là một phần của một con
số cực lớn. Trả lời được nghĩa là thật sự nắm giữ bí mật.

\section{AnonCreds là gì}
\label{sec:anoncreds}

AnonCreds (\textit{Anonymous Credentials}) là bộ tiêu chuẩn mô tả cách áp dụng ý tưởng
trên vào giấy tờ định danh thật. Nó được dùng trong các hệ thống định danh quốc gia như
Bhutan NDI. NX Cred hiện thực lại phần lõi của tiêu chuẩn này.

AnonCreds đóng góp bốn ý tưởng, và đây là bốn ý tưởng nên nhớ nếu chỉ nhớ một mục:

\begin{description}[leftmargin=0pt,style=nextline,itemsep=6pt]
  \item[Link secret --- bí mật gắn thẻ với chủ thẻ]
    Ví tự sinh ra một con số ngẫu nhiên và không bao giờ cho ai xem, kể cả cơ quan cấp
    thẻ. Thẻ được ký sao cho chỉ dùng được cùng với con số đó. Hệ quả: ai đó chép trộm
    được file thẻ của bạn cũng không dùng được, vì họ không có link secret.

  \item[Ký mù --- cơ quan cấp ký lên thứ nó không đọc được]
    Ví gửi lên một \emph{cam kết} che kín link secret. Cơ quan cấp ký lên cam kết đó.
    Chữ ký ràng buộc chặt vào link secret, nhưng cơ quan cấp chưa từng nhìn thấy nó.

  \item[Tiết lộ chọn lọc --- chỉ mở đúng ô được hỏi]
    Thẻ có tám trường. Bên xác minh hỏi một trường thì bảy trường còn lại vẫn nằm trong
    bằng chứng dưới dạng đã che, đủ để chứng minh thẻ hợp lệ nhưng không lộ giá trị.

  \item[Không liên kết được --- mỗi lần trình là một lần mới]
    Mỗi lần trình thẻ, ví ngẫu nhiên hoá lại chữ ký. Hai bên xác minh có gộp dữ liệu với
    nhau cũng không tìm ra điểm chung nào để kết luận ``hai lần này là cùng một người''.
\end{description}

\section{NX Cred là gì}
\label{sec:nxcred}

NX Cred là hiện thực của mô hình trên: một ví định danh chạy trên trình duyệt, một cơ quan
cấp thẻ mô phỏng, một cổng dành cho bên xác minh, và một hợp đồng trên blockchain giữ khoá
công khai.

Giấy tờ được dùng là \textbf{căn cước công dân}, với đúng tám trường in ở mặt trước:

\begin{center}
\begin{tabular}{llll}
\toprule
\sffamily Số / No. & \sffamily Họ và tên & \sffamily Ngày sinh & \sffamily Giới tính \\
\sffamily Quốc tịch & \sffamily Quê quán & \sffamily Nơi thường trú & \sffamily Có giá trị đến \\
\bottomrule
\end{tabular}
\end{center}

Cơ quan cấp không ký bất cứ thứ gì ngoài tám trường này, và bên xác minh cũng không hỏi
được gì ngoài tám trường này. Ràng buộc đó được thi hành bằng khuôn mẫu đã đăng trên chain,
không phải bằng thoả thuận miệng.

Một điểm khác biệt về quyền lực đáng chú ý: trong mô hình truyền thống, mỗi lần bạn xác
minh danh tính là một lần dữ liệu của bạn được nhân bản thêm vào một cơ sở dữ liệu mới.
Mười dịch vụ là mười bản sao, và chỉ cần một trong số đó bị tấn công là dữ liệu của bạn ra
ngoài. Với NX Cred, mười dịch vụ nhận về mười bằng chứng toán học khác nhau, không bản nào
chứa dữ liệu gốc, và không bản nào liên kết được với bản còn lại.

\vspace{4mm}
\hrule
\vspace{6mm}

\part*{\sffamily NX stack}
\addcontentsline{toc}{part}{NX stack}

\section{Issuer --- cơ quan cấp thẻ}

Issuer là bên có thẩm quyền xác nhận danh tính: trong thực tế là cơ quan quản lý căn cước,
trong bản demo là một dịch vụ mô phỏng có sẵn một cơ sở dữ liệu hồ sơ.

Issuer làm bốn việc:

\begin{enumerate}[leftmargin=1.4em,itemsep=4pt]
  \item \textbf{Sinh cặp khoá} và đăng phần công khai lên chain, kèm khuôn mẫu tám trường.
  \item \textbf{Đối chiếu eKYC} --- kiểm tra dữ liệu người dùng khai có khớp hồ sơ hay không.
  \item \textbf{Kiểm bằng chứng nhỏ} mà ví gửi kèm, để chắc ví thật sự nắm link secret bên
        trong cam kết chứ không phải nhặt cam kết của người khác.
  \item \textbf{Ký mù} và trả thẻ đã ký về cho ví.
\end{enumerate}

Issuer \emph{không} lưu thẻ đã cấp, \emph{không} biết thẻ được dùng ở đâu, và \emph{không}
bao giờ thấy link secret. Nó chỉ ghi nhận một việc duy nhất: số căn cước này đã được cấp
thẻ rồi. Nhờ vậy một số căn cước không thể có hai ví cùng lúc.

\section{Holder --- ví của người dùng}

Holder là phần mềm chạy phía người dùng, giữ ba thứ và cả ba đều không được rời khỏi thiết
bị trong mô hình đầy đủ:

\begin{center}
\begin{tabular}{@{}p{3.4cm}p{10.4cm}@{}}
\toprule
\textbf{\sffamily Link secret} & Con số ngẫu nhiên gắn thẻ với chính ví này. Mất nó thì thẻ
thành vô dụng; lộ nó thì người khác dùng được thẻ của bạn. \\
\addlinespace
\textbf{\sffamily Credential} & Chữ ký của cơ quan cấp lên toàn bộ tám thuộc tính cộng với
link secret. \\
\addlinespace
\textbf{\sffamily Passkey} & Khoá sinh trắc học (vân tay, khuôn mặt) mở kho lưu trữ. Đây là
lớp bảo vệ quyền truy cập, không phải cơ chế mật mã của bằng chứng. \\
\bottomrule
\end{tabular}
\end{center}

Ví là bên duy nhất có thể tạo ra bằng chứng, vì chỉ nó biết link secret. Mỗi lần trình
thẻ, ví dựng một bằng chứng hoàn toàn mới.

\section{Verifier --- bên xác minh}

Verifier là bất kỳ ai cần kiểm tra danh tính: ngân hàng, khách sạn, nhà tuyển dụng.

Verifier soạn một \emph{yêu cầu trình bày} gồm hai phần: những trường cần được tiết lộ, và
những điều cần được chứng minh mà không tiết lộ. Kèm theo là một số ngẫu nhiên dùng một lần
để bằng chứng cũ không thể phát lại.

Điều quan trọng về mặt tin cậy: verifier lấy khoá công khai \emph{từ chain}, không phải từ
người đang trình thẻ. Nếu nhận khoá từ người trình, kẻ gian chỉ việc tự tạo khoá riêng, tự
ký một tấm thẻ giả rồi đưa kèm khoá công khai của chính mình --- mọi phép kiểm đều qua.
Lấy khoá từ nguồn độc lập là thứ chặn đường tấn công đó.

\vspace{4mm}
\hrule
\vspace{6mm}

\part*{\sffamily Nền tảng thiết kế}
\addcontentsline{toc}{part}{Nền tảng thiết kế}

\section{Luồng giữa Issuer và Holder --- cấp thẻ}
\label{sec:flow-ih}

Đây là luồng chạy \textbf{một lần duy nhất} cho mỗi người dùng.

\begin{center}
\begin{tikzpicture}[node distance=9mm]
  \node[actor] (h1) {Holder};
  \node[actor,right=68mm of h1] (i1) {Issuer};

  \node[blk,below=8mm of h1,text width=44mm] (a) {1. Chụp căn cước, trích xuất tám trường};
  \node[blk,below=8mm of i1,text width=44mm] (b) {2. Đối chiếu hồ sơ eKYC\\ phát số dùng một lần};

  \node[blk,below=8mm of a,text width=44mm] (c) {3. Sinh \textbf{link secret}\\ tính cam kết che kín nó};
  \node[blk,below=8mm of c,text width=44mm] (d) {4. Kèm bằng chứng nhỏ:\\ ``tôi biết cái bên trong cam kết''};

  \node[blk,below=22mm of b,text width=44mm] (e) {5. Kiểm bằng chứng nhỏ};
  \node[blk,below=8mm of e,text width=44mm] (f) {6. \textbf{Ký mù} lên cam kết\\ và tám thuộc tính};

  \node[blk,below=8mm of d,text width=44mm,yshift=-8mm] (g) {7. Gỡ mù, tự kiểm chữ ký\\ rồi lưu thẻ vào ví};

  \draw[step] (a.east) -- node[above]{dữ liệu eKYC} (b.west);
  \draw[step] (b.west) to[out=200,in=340] node[below,pos=0.5]{số dùng một lần} (c.east);
  \draw[step] (d.east) -- node[above]{cam kết + bằng chứng} (e.west);
  \draw[step] (e) -- (f);
  \draw[step] (f.west) to[out=200,in=20] node[above,pos=0.55]{thẻ đã ký} (g.east);

  \node[note,below=6mm of g,text width=110mm]
    {Link secret chưa bao giờ rời khỏi ví. Issuer ký lên một thứ nó không đọc được.};
\end{tikzpicture}
\end{center}

Bước 4 là bước dễ bị bỏ qua nhưng không thể thiếu. Nếu chỉ gửi cam kết mà không kèm bằng
chứng, một kẻ gian có thể chặn lấy cam kết của người khác rồi nộp như của mình. Bằng chứng
nhỏ này buộc người nộp phải thật sự biết con số nằm bên trong.

Bước 7 cũng vậy: ví tự kiểm chữ ký trước khi lưu, để cơ quan cấp không thể trả về một chữ
ký hỏng rồi sau này đổ lỗi cho ví.

\section{Luồng giữa Holder và Verifier --- xác minh}
\label{sec:flow-hv}

Đây là luồng chạy \textbf{mỗi lần} bạn cần chứng minh danh tính.

\begin{center}
\begin{tikzpicture}[node distance=9mm]
  \node[actor] (v1) {Verifier};
  \node[actor,right=64mm of v1] (h1) {Holder};

  \node[blk,below=8mm of v1,text width=46mm] (a) {1. Soạn yêu cầu:\\ cần lộ trường nào, cần chứng minh điều gì\\ kèm số ngẫu nhiên dùng một lần};

  \node[blk,below=8mm of h1,text width=46mm] (b) {2. Người dùng xem yêu cầu\\ và \textbf{tự quyết định} chia sẻ gì};
  \node[blk,below=8mm of b,text width=46mm] (c) {3. Ngẫu nhiên hoá lại chữ ký};
  \node[blk,below=8mm of c,text width=46mm] (d) {4. Che link secret và các trường không chia sẻ};
  \node[blk,below=8mm of d,text width=46mm] (e) {5. Tạo bằng chứng};

  \node[blk,below=44mm of a,text width=46mm] (f) {6. Lấy khoá công khai \textbf{từ chain}};
  \node[blk,below=8mm of f,text width=46mm] (g) {7. Kiểm bằng chứng\\ hợp lệ hay không};

  \draw[step] (a.east) -- node[above]{yêu cầu} (b.west);
  \draw[step] (b) -- (c);
  \draw[step] (c) -- (d);
  \draw[step] (d) -- (e);
  \draw[step] (e.west) to[out=200,in=340] node[below,pos=0.5]{bằng chứng} (f.east);
  \draw[step] (f) -- (g);

  \node[note,below=6mm of g,text width=120mm]
    {Verifier kết thúc với đúng một chữ ``hợp lệ'', cộng với những trường bạn đồng ý mở.
     Không có bản sao dữ liệu gốc nào được tạo ra.};
\end{tikzpicture}
\end{center}

Bước 3 là thứ tạo ra tính không liên kết được. Chữ ký gửi đi mỗi lần một khác, nên hai bên
xác minh không thể đối chiếu sổ sách để lần ra bạn.

Bước 2 đáng nhấn mạnh về mặt sản phẩm: quyền quyết định nằm ở người dùng, không nằm ở bên
xác minh. Bên xác minh chỉ được \emph{hỏi}; bên trả lời là ví, và ví hỏi chủ của nó trước.

\section{Luồng giữa Issuer và Verifier --- không có luồng nào}
\label{sec:flow-iv}

Đây là mục quan trọng nhất của tài liệu, và nội dung của nó là một sự vắng mặt.

\begin{center}
\begin{tikzpicture}[node distance=14mm]
  \node[actor] (i) {Issuer};
  \node[chain,right=34mm of i] (c) {Chain};
  \node[actor,right=34mm of c] (v) {Verifier};

  \draw[step] (i) -- node[above,align=center]{đăng một lần\\\scriptsize khoá công khai, khuôn mẫu} (c);
  \draw[step] (c) -- node[above,align=center]{đọc khi cần\\\scriptsize không cần xin phép} (v);

  \draw[-{Stealth[length=2.4mm]},dashed,draw=muted,thick]
    (i) to[out=-30,in=210] node[below,note,yshift=-1mm]{không có kênh nào ở đây} (v);
\end{tikzpicture}
\end{center}

Trong mô hình cũ, mỗi lần bạn dùng giấy tờ ở đâu đó, bên đó phải hỏi lại cơ quan cấp:
``thẻ này có thật không?''. Câu hỏi đó biến cơ quan cấp thành một nơi biết toàn bộ lịch sử
đi lại, vay mượn, thuê nhà của bạn --- một dạng giám sát phát sinh ngoài ý muốn.

NX Cred cắt đứt kênh đó. Cơ quan cấp đăng khoá công khai lên chain \emph{một lần}, rồi
không tham gia gì nữa. Bên xác minh tự đọc, tự kiểm. Cơ quan cấp không biết, và về mặt kỹ
thuật là \emph{không thể} biết, thẻ của bạn đang được dùng ở đâu.

\section{Blockchain được dùng ở đâu}

Blockchain trong hệ thống này có một vai trò hẹp và rõ ràng. Nó \textbf{không} lưu dữ liệu
cá nhân, \textbf{không} lưu thẻ, và \textbf{không} ghi lại các lần xác minh.

\begin{center}
\begin{tabular}{@{}p{6.6cm}p{7.4cm}@{}}
\toprule
\textbf{\sffamily Có trên chain} & \textbf{\sffamily Không lên chain} \\
\midrule
Khoá công khai của cơ quan cấp & Dữ liệu cá nhân của bất kỳ ai \\
Khuôn mẫu tám trường của giấy tờ & Thẻ đã cấp \\
Danh sách cơ quan cấp được tin cậy & Link secret \\
 & Lịch sử ai xác minh cái gì, ở đâu \\
\bottomrule
\end{tabular}
\end{center}

Lý do cần đến blockchain chỉ có một: bên xác minh phải lấy được khoá công khai từ một nguồn
mà \emph{không ai sửa được}, kể cả chính cơ quan cấp. Nếu khoá nằm trên một máy chủ thông
thường, người kiểm soát máy chủ đó có thể âm thầm đổi khoá và từ đó ký được thẻ giả mà
không ai phát hiện. Sổ cái công khai loại bỏ khả năng đó: mọi thay đổi đều để lại dấu vết
và ai cũng kiểm tra được.

Nói cách khác, blockchain ở đây đóng vai một \emph{bảng tin công khai không thể tẩy xoá},
chứ không phải một cơ sở dữ liệu.

\vspace{4mm}
\hrule
\vspace{6mm}

\part*{\sffamily Nâng cao}
\addcontentsline{toc}{part}{Nâng cao}

\section{Lõi mật mã (đọc thêm)}
\label{sec:crypto}

Toàn bộ phần toán học --- công thức sinh khoá, cam kết, giao thức sigma, chữ ký mù, bằng
chứng trình bày, và chứng minh tính đúng đắn của từng phép kiểm --- nằm trong một tài liệu
riêng:

\begin{center}
\fbox{\begin{minipage}{0.86\textwidth}
\vspace{2mm}
\centering
{\sffamily\large\bfseries Lõi mật mã NX Cred}\\[3pt]
{\color{muted}\small Chữ ký Camenisch--Lysyanskaya và bằng chứng không tiết lộ ---
đầy đủ định lý và chứng minh}\\[6pt]
\href{run:./nxcred-crypto.pdf}{\texttt{nxcred-crypto.pdf}}
\vspace{2mm}
\end{minipage}}
\end{center}

\noindent Tài liệu đó dành cho người muốn thẩm định, không cần thiết để hiểu sản phẩm. Nội
dung gồm: cấu trúc nhóm và lý do chọn safe prime, tính giấu kín hoàn hảo của cam kết, tính
đầy đủ và tính chắc chắn của giao thức sigma, chứng minh đẳng thức chữ ký được bảo toàn qua
phép ngẫu nhiên hoá, và chứng minh bên xác minh khôi phục lại đúng giá trị cần so sánh.

\section{Các thư viện sử dụng}

\subsection*{Phía máy chủ --- Python}

\begin{longtable}{@{}p{2.9cm}p{6.6cm}p{4.4cm}@{}}
\toprule
\textbf{\sffamily Thư viện} & \textbf{\sffamily Dùng để làm gì} & \textbf{\sffamily Tải về} \\
\midrule
\endhead
gmpy2 & Số học số lớn: sinh safe prime, luỹ thừa modulo. Đây là thư viện gánh toàn bộ phần nặng của lõi mật mã. & \href{https://pypi.org/project/gmpy2/}{pypi.org/project/gmpy2} \\
\addlinespace
FastAPI & Khung web dựng các endpoint của issuer, ví và verifier. & \href{https://fastapi.tiangolo.com/}{fastapi.tiangolo.com} \\
\addlinespace
Uvicorn & Máy chủ chạy ứng dụng FastAPI. & \href{https://www.uvicorn.org/}{uvicorn.org} \\
\addlinespace
Pydantic & Kiểm tra và ràng buộc kiểu dữ liệu vào ra của API. & \href{https://pypi.org/project/pydantic/}{pypi.org/project/pydantic} \\
\addlinespace
py\_webauthn & Xác minh passkey theo chuẩn WebAuthn: phát thách thức, kiểm chữ ký thiết bị. & \href{https://pypi.org/project/webauthn/}{pypi.org/project/webauthn} \\
\addlinespace
google-auth & Kiểm tra ID-token khi người dùng đăng nhập bằng tài khoản Google. & \href{https://pypi.org/project/google-auth/}{pypi.org/project/google-auth} \\
\addlinespace
web3.py & Đọc khoá công khai và khuôn mẫu giấy tờ từ hợp đồng trên chain. & \href{https://pypi.org/project/web3/}{pypi.org/project/web3} \\
\addlinespace
websockets & Kênh thời gian thực giữa máy tính và điện thoại khi quét giấy tờ. & \href{https://pypi.org/project/websockets/}{pypi.org/project/websockets} \\
\bottomrule
\end{longtable}

\subsection*{Phía trình duyệt --- TypeScript}

\begin{longtable}{@{}p{2.9cm}p{6.6cm}p{4.4cm}@{}}
\toprule
\textbf{\sffamily Thư viện} & \textbf{\sffamily Dùng để làm gì} & \textbf{\sffamily Tải về} \\
\midrule
\endhead
React & Dựng giao diện ví và cổng xác minh. & \href{https://react.dev/}{react.dev} \\
\addlinespace
Vite & Công cụ đóng gói và máy chủ phát triển. & \href{https://vite.dev/}{vite.dev} \\
\addlinespace
TypeScript & Ngôn ngữ có kiểu tĩnh, bắt lỗi trước khi chạy. & \href{https://www.typescriptlang.org/}{typescriptlang.org} \\
\addlinespace
Tailwind CSS & Hệ thống tạo kiểu giao diện. & \href{https://tailwindcss.com/}{tailwindcss.com} \\
\addlinespace
Framer Motion & Hiệu ứng chuyển cảnh giữa các bước. & \href{https://motion.dev/}{motion.dev} \\
\addlinespace
jsQR & Đọc mã QR in trên mặt trước căn cước. & \href{https://www.npmjs.com/package/jsqr}{npmjs.com/package/jsqr} \\
\addlinespace
Tesseract.js & Đọc chữ trên ảnh giấy tờ, bổ sung cho các trường mã QR không chứa. & \href{https://tesseract.projectnaptha.com/}{tesseract.projectnaptha.com} \\
\addlinespace
qrcode & Sinh mã QR để chuyển phiên từ máy tính sang điện thoại. & \href{https://www.npmjs.com/package/qrcode}{npmjs.com/package/qrcode} \\
\addlinespace
React Router & Điều hướng giữa các màn hình. & \href{https://reactrouter.com/}{reactrouter.com} \\
\bottomrule
\end{longtable}

\vspace{4mm}
\hrule
\vspace{6mm}

\part*{\sffamily Hướng dẫn sử dụng}
\addcontentsline{toc}{part}{Hướng dẫn sử dụng}

\section{Hướng dẫn sử dụng}

\vspace{4mm}
{\color{muted}\itshape Phần này sẽ được bổ sung sau.}

\end{document}
